\begin{abstract}
Multi-agent LLM orchestration is increasingly framed as a routing problem, an aggregation problem, or a post-hoc failure problem. We study an intermediate object that none of these frames opens up: the responsibility structure an agent projects between receiving a natural-language delegation and executing against an artifact. We formalize \emph{responsibility projection} as multi-label weight prediction over a closed dimension set, instantiate it on $J_{v1.1}$ --- a 12-dimension taxonomy for paper-research delegation (seven category dimensions, five cross-cutting) --- and use the closure to make projections from different model families directly comparable.

The primary empirical contribution (P1) is that pre-execution responsibility projection is measurable and family-attributable. On a 50-example pilot, cross-family projection mismatch under three distinct model families is approximately $17\times$ within-family stochastic variance under five-run repetition (median $R(d) = 14.4$, paired bootstrap CI excludes zero, Wilcoxon $p < 0.0001$). The secondary contribution (P2) is a negative actionability result: under a 12-judge Anthropic-excluded panel and a three-condition execution split, projection-driven execution shows a directional \emph{disadvantage} relative to direct execution on the headline weighted-R1 settlement loss (loss \emph{increase} for projection\_driven over direct\_naive of $0.139$ on the mean; equivalently, paired diff direct\_naive $-$ projection\_driven $= -0.139$ with 95\% bootstrap CI $[-0.169, -0.109]$). The cost concentrates on R1.4 (novelty mapping) and R1.7 (citation audit) --- the two dimensions whose $s = 5$ anchor demands deep specialty engagement. We report this as a boundary condition, not a refutation of the projection layer: responsibility projection is useful as a measurement and governance object, but naive injection of the projected vector into the executor can dilute deep dimension-specific work unless paired with retrieval, comparator corpora, or domain-expert anchoring. The methodological contribution (P3) is a closed Stage 1 human-anchor pilot on R1.7 that surfaces an \emph{anchor specifiability ceiling}: form-embedded sharpened anchors are insufficient for reliable rater application without a separately-read protocol document, and the human-anchor scalability constraint is anchor specifiability and domain-expertise gating, not rater throughput.

The scope of this paper is one delegation category (R1, paper-research) and a pilot sample ($n = 50$). The full closed-loop protocol (bid, clarification, selection, settlement, reputation) is formalized but evaluated only on the projection layer; the full four-condition Experiment 2 with task-aware-routing and CLAMBER-style baselines, longitudinal reputation, and main-run sample expansion are reserved for follow-up work.
\end{abstract}

\begin{IEEEkeywords}
Pre-execution responsibility projection, multi-agent LLM orchestration, closed responsibility taxonomy, projection mismatch, settlement loss, anchor specifiability, controlled evaluation.
\end{IEEEkeywords}
